\documentclass{article}
\usepackage{/home/user1/code/tex/sty}
\usepackage{amsfonts}
\usepackage{amsmath}
\usepackage{geometry}
\usepackage{graphicx}
\usepackage{hyperref}
\usepackage{floatrow}
\newcommand{\fig}[3]{
\begin{figure}[H]
\begin{center}
\includegraphics[height=#1]{#2}
\caption{#3}
\label{#2}
\end{center}
\end{figure}}
\usepackage[backend=bibtex]{biblatex}
\addbibresource{bib.bib}
\setlength{\parindent}{0pt}
\geometry{top=1in,bottom=1in,left=1in,right=1in}
\begin{document}
\begin{center}
\textbf{\Large{Problem Set 2 Corrections}}\\~\\
\begin{tabular}{l l}
Student&Agustin Romero\\
Email&ar1@stanford.edu\\
Instructor&Caterina Vernieri\\
Course&Physics 252\\
Institution&Stanford University\\
Date&\today\\
\end{tabular}
\end{center}
\section*{Problem 1}
The student answer was not correct. This problem asked to explain which particles in the Standard Model are irrelevant to a typical chemistry laboratory. The student's answer listed all quarks, leptons, and gauge bosons, and explained which ones are not relevant, but neglected to explain why neutrinos and antiparticles are irrelevant to the typical chemistry laboratory.\\~\\
To correct the solution, include an explanation for why almost all antiparticles and neutrinos are irrelevant to chemistry. $\mu$, $\tau$, $\nu_\mu$ and $\nu_\tau$ and their antiparticles are almost always irrelevant to chemistry because of their low luminosity in the laboratory. However $\nu_\mu$ and $\bar{\nu}_\mu$ are relevant to chemistry from a theoretical perspective, in the $\beta$ decay process, even though the resulting neutrino will very likely not interact in the laboratory.
\section*{Problem 2.a.}
The student answer was correct.
\section*{Problem 2.b.}
The student answer was correct.
\section*{Problem 2.c.}
The student answer was correct.
\section*{Problem 2.d.}
The student answer was not correct. The student answer stated the process was not Cabibbo suppressed, but it in fact is. To fix the answer, state the \fbox{process is Cabbibo suppressed}.
\section*{Problem 2.e.}
The student answer was not correct. The question asked to answer if a process was allowed and what its diagrams were. The student answer neglected to check if the process was kinematically allowed.
\e{m_p=938\un{MeV}}
\e{m_{K^-}=493\un{MeV}}
\e{m_{\Sigma^0}=1192\un{MeV}}
To fix the solution, state the process is \fbox{not kinematically allowed}.
\section*{Problem 2.f.}
The student answer was correct.
\section*{Problem 2.g.}
The student answer was correct.
\section*{Problem 2.h.}
The student answer was correct.
\section*{Problem 2.i.}
The student answer was not correct. The student answer drew a diagram for the process, but instead, the process it not allowed by charge conservation. To fix the solution, replace the solution with the statement that \fbox{conservation of charge} is violated.
\section*{Problem 3.a.}
The student answer was correct.
\section*{Problem 3.b.}
The student answer was correct.
\section*{Problem 4.}
The student answer was not correct. The student answer started with the invariance of $s$, which could lead to the correct answer, except that $\vec{p_\gamma}=-\vec{p_p}$ was asserted, which is not correct. To fix the solution, start with the \fbox{invariance of $p_\mu p^\mu$ in the initial and final state}.
\section*{Problem 5.a.}
The student answer was correct.
\section*{Problem 5.b.}
The student answer was correct.
\section*{Problem 5.c.}
The student answer was correct. The correct calculation is:\\~\\
Label the process $p+p\rightarrow p + p + \bar{p} + \bar{p}$ as $1+2\rightarrow 3+4+5+6$. One initial proton is at rest, and all final state protons are at rest. So let all $\vec{p}=0$ except one of the initial state protons, $\vec{p_2}$.
\e{\vec{p_1}=\vec{p_3}=\vec{p_4}=\vec{p_5}=\vec{p_6}=0}
Then the energy of particle 1 is
\e{E_1=\gamma_1 m_1 c^2 = m_1 c^2 = m_p c^2}
Use the invariance of $s$ to solve for the energy of the energy of the moving proton, particle 2.
\e{(\sum E_i)^2-(\sum \vec{p_i})^2 &= (\sum E_j)^2-(\sum \vec{p_j})^2\\
(E_1+E_2)^2 + \vec{p_2}^2 &= 16m_p^2 c^4 \\
E_1^2 + E_1E_2 + E_2^2 - p_2^2 &= 16m_p^2 c^4\\
2E_1E_2+2m_2^2c_4=16m_p^2c_4\\
2E_1E_2=14m_p^2c_4\\
2m_pc^2E_2=14m_p^2c^4\\
\boxed{E_2=7m_pc^2}}
In agreement with the solution.
\section*{Problem 6.a.}
The student answer was correct.
\section*{Problem 6.b.}
The student answer was not correct. The diagram for $\Delta^+ \rightarrow n\pi^+$ used a photon propagator to create the pion, which is wrong. Use a gluon propagator like so.
\fig{5cm}{1.png}{Fixed diagram for the first process.}
\section*{Problem 6.c.}
The student answer was not correct. The student answer stated the process for $\Delta^+ \rightarrow n\pi^+$ is electromagnetic, but this is wrong. To fix this solution, state the first process is \fbox{a strong interaction}, as corrected in 6.b. 
\end{document}
